El \textit{World Editor} es una herramienta destinada a la creación y edición de mundos dentro del juego.
Está construido como un cliente de Unity que se comunica con el \textit{Core Service} a través del \textit{API Gateway}.
El editor permite a los jugadores crear nuevos mundos,
modificar la configuración de zonas existentes, gestionar \textit{assets}, publicar y mediante una subscripción activar dichos mundos.

\subsubsection{Importación de modelos 3D}
Una de las funcionalidades centrales del editor es la capacidad de importar modelos 3D
personalizados por los creadores. Sin embargo, Unity no permite instanciar modelos 3D
en tiempo de ejecución a partir de archivos externos, ya que sus mecanismos de carga de
recursos están pensados para activos conocidos en tiempo de compilación. Para resolver
esta limitación se optó por la librería GLTFast, que permite importar modelos en formato
\texttt{.glb} de forma dinámica durante la ejecución.

\vspace{0.25cm}

El formato \texttt{.glb} es un formato abierto y ampliamente soportado para la
representación de escenas 3D, con la ventaja de contener tanto la geometría del modelo
como sus texturas y materiales en un único archivo, lo que simplifica el proceso de
importación y la gestión de recursos. GLTFast se utiliza no solo en el editor de mundos,
sino también en el cliente del juego, que descarga los modelos desde el Core Service.


\subsubsection{Gestión de Mundos}

El editor cuenta con un sistema de guardado y apertura de mundos similar al de otros programas de edición,
permitiendo retomar la edición en distintas sesiones de trabajo.

Los mundos se persisten localmente utilizando \texttt{Application.persistentDataPath} de Unity, una ruta
del sistema de archivos gestionada por el motor que garantiza permisos de escritura en todas las plataformas
soportadas. Cada mundo se almacena como un directorio independiente con el nombre del mundo, con la
siguiente estructura:

\begin{verbatim}
<mundo>/
|- world.data
|- creatables.data
|- zone_1.world
|- zone_2.world
|- ...
|- zone_n.world
\end{verbatim}

\begin{itemize}
  \item \texttt{world.data}: Contiene la configuración general del mundo, incluyendo nombre, descripción
  e identificador asignado al momento de su publicación.
  \item \texttt{creatables.data}: Almacena la definición de todos los objetos creables configurados
  para el mundo, como \textit{NPCs}, enemigos, ítems y misiones.
  \item \texttt{zone\_\{n\}.world}: Representa el estado de una zona independiente, almacenando todos
  los \textit{Placeables} colocados en ella junto con su posición, rotación, escala y propiedades configuradas.
  Se genera un archivo por cada zona que compone el mundo.
\end{itemize}


\subsubsection{Gestión de Recursos en Memoria}

La gestión de recursos en memoria se organiza mediante un sistema de managers y repositorios internos, 
donde cada tipo de dato cuenta con su propio repositorio especializado que encapsula la lógica de 
serialización y acceso a disco de su dominio.

Los \textit{Placeables}, al estar representados físicamente en la escena, son rastreados individualmente y 
notifican al sistema mediante eventos cuando su estado es modificado, determinando si deben incluirse 
en el próximo ciclo de guardado.

Los \textit{Creatables}, al no tener representación física en la escena, son gestionados por un manager dedicado 
que mantiene en memoria el conjunto de objetos creables del mundo, permitiendo su consulta, modificación 
y persistencia de forma centralizada.

\subsubsection{Publicación de Mundos}

Una vez finalizado el diseño, el creador puede publicar su mundo para que pueda ser accedido mediante el cliente de juego.
El proceso requiere haber iniciado sesión y se realiza desde el panel de publicación. Al confirmar la publicación,
se inicia una secuencia de carga hacia el \textit{Core Service} que sigue el siguiente orden:

\begin{enumerate}
  \item \textbf{Datos del Mundo (\textit{World Data})}: Los datos estructurales del mundo en formato \textit{JSON}, que incluyen la configuración
  de zonas, objetos y sus propiedades.
  \item \textbf{Modelos 3D}: Los modelos utilizados en el mundo, excluyendo aquellos que forman parte del
  conjunto de \textit{assets} predeterminados del sistema.
  \item \textbf{\textit{Assets} de Ítems}: Las imágenes e iconos asociados a los ítems creados por el usuario.
  \item \textbf{Cosméticos}: Los \textit{assets} visuales correspondientes a los cosméticos configurados en el mundo.
  \item \textbf{\textit{Assets} de Entorno}: Las texturas de suelo (\textit{ground textures}) y cielos
  (\textit{skybox textures}) asociadas a cada zona del mundo.
\end{enumerate}

Luego que el mundo es publicado, este debe ser activado mediante el panel de economía, lo cual requiere una cierta cantidad de \textit{slots}
igual a la cantidad de zonas que se quieran desplegar y estos se obtienen a partir de suscripciones.

\vspace{0.25cm}

Una optimización importante del proceso de publicación es evitar la recarga de recursos
que ya se encuentran en el servidor. Esto se logra omitiendo tanto los recursos que
forman parte del conjunto de \textit{assets} predeterminados del sistema como los modelos
3D personalizados que ya hayan sido subidos en publicaciones anteriores.
